\documentclass[12pt]{article}
\usepackage[a4paper,margin=1in]{geometry}
\usepackage{amsmath,amssymb}
\usepackage{mathtools}
\usepackage{bm}

\title{Positive-$P$ Derivation for the Open Bose-Hubbard Model}
\author{}
\date{}

\begin{document}
\maketitle

\section*{Purpose}
This note derives the positive-$P$ equations for the open Bose-Hubbard model with onsite interaction strength $U$, single-particle loss rate $\gamma$, and nearest-neighbor hopping amplitude $J$. The derivation follows the same logic as the outline given in Appendix A of PRX Quantum \textbf{2}, 010319 (2021), but it is written in thesis-friendly notation and then specialized to the single-site stochastic differential equations (SDEs) currently used in this project.

\section*{Model}
Consider bosonic annihilation operators $\hat a_j$ on lattice sites $j=1,\dots,M$. The Hamiltonian is
\[
\hat H
=
\sum_{j=1}^M \frac{U}{2}\,\hat a_j^{\dagger 2}\hat a_j^2
-J\sum_{\langle j,k\rangle}
\left(
\hat a_j^\dagger \hat a_k + \hat a_k^\dagger \hat a_j
\right),
\]
and the local loss operators are
\[
\hat L_j = \sqrt{\gamma}\,\hat a_j.
\]
The density operator obeys the Lindblad master equation
\[
\frac{d\hat\rho}{dt}
=
-i[\hat H,\hat\rho]
+
\gamma \sum_{j=1}^M
\left(
\hat a_j \hat\rho \hat a_j^\dagger
- \frac12 \hat a_j^\dagger \hat a_j \hat\rho
- \frac12 \hat\rho \hat a_j^\dagger \hat a_j
\right).
\]

\section*{Positive-$P$ Representation}
In the positive-$P$ representation, the density operator is expanded as
\[
\hat\rho
=
\int P(\bm{\alpha},\bm{\widetilde{\alpha}})\,
\hat\Lambda(\bm{\alpha},\bm{\widetilde{\alpha}})\,
d^{2M}\bm{\alpha}\,d^{2M}\bm{\widetilde{\alpha}},
\]
with kernel
\[
\hat\Lambda(\bm{\alpha},\bm{\widetilde{\alpha}})
=
\bigotimes_{j=1}^M
\frac{|\alpha_j\rangle\langle \widetilde{\alpha}_j^*|}
{\langle \widetilde{\alpha}_j^*|\alpha_j\rangle}.
\]
Here $\alpha_j$ and $\widetilde{\alpha}_j$ are independent complex variables. In general one must \emph{not} set $\widetilde{\alpha}_j = \alpha_j^*$ trajectory by trajectory. Physical observables are obtained only after stochastic averaging.

For each site, the kernel satisfies the differential identities
\begin{align}
\hat a_j \hat\Lambda &= \alpha_j \hat\Lambda, \\
\hat a_j^\dagger \hat\Lambda &= \left(\widetilde{\alpha}_j + \frac{\partial}{\partial \alpha_j}\right)\hat\Lambda, \\
\hat\Lambda \hat a_j^\dagger &= \widetilde{\alpha}_j \hat\Lambda, \\
\hat\Lambda \hat a_j &= \left(\alpha_j + \frac{\partial}{\partial \widetilde{\alpha}_j}\right)\hat\Lambda.
\end{align}
These identities allow every operator product in the master equation to be converted into differential operators acting on $\hat\Lambda$.

\section*{Contribution of the Interaction Term}
For the onsite interaction at site $j$,
\[
\hat H_{U,j} = \frac{U}{2}\hat a_j^{\dagger 2}\hat a_j^2,
\]
the commutator contribution is
\[
-i[\hat H_{U,j},\hat\rho].
\]
After inserting the positive-$P$ expansion and using the kernel identities, one obtains the standard second-order differential structure
\[
\frac{\partial P}{\partial t}\Big|_{U,j}
=
-\frac{\partial}{\partial \alpha_j}
\left(
-iU\,\widetilde{\alpha}_j\alpha_j^2\,P
\right)
-\frac{\partial}{\partial \widetilde{\alpha}_j}
\left(
+iU\,\widetilde{\alpha}_j^{\,2}\alpha_j\,P
\right)
+ \frac12 \frac{\partial^2}{\partial \alpha_j^2}
\left(
-iU\,\alpha_j^2\,P
\right)
+ \frac12 \frac{\partial^2}{\partial \widetilde{\alpha}_j^{\,2}}
\left(
+iU\,\widetilde{\alpha}_j^{\,2}\,P
\right).
\]
The first two terms are the nonlinear drift generated by the Kerr interaction, while the second-order terms give the multiplicative noise.

\section*{Contribution of the Dissipation Term}
For local single-particle loss,
\[
\gamma
\left(
\hat a_j \hat\rho \hat a_j^\dagger
- \frac12 \hat a_j^\dagger \hat a_j \hat\rho
- \frac12 \hat\rho \hat a_j^\dagger \hat a_j
\right),
\]
the positive-$P$ identities give a purely first-order contribution,
\[
\frac{\partial P}{\partial t}\Big|_{\gamma,j}
=
-\frac{\partial}{\partial \alpha_j}
\left(
-\frac{\gamma}{2}\alpha_j\,P
\right)
-\frac{\partial}{\partial \widetilde{\alpha}_j}
\left(
-\frac{\gamma}{2}\widetilde{\alpha}_j\,P
\right).
\]
Thus, dissipation contributes only linear damping to the drift and no diffusion term.

\section*{Contribution of the Hopping Term}
For nearest-neighbor hopping,
\[
\hat H_J
=
-J\sum_{\langle j,k\rangle}
\left(
\hat a_j^\dagger \hat a_k + \hat a_k^\dagger \hat a_j
\right),
\]
the master-equation contribution is again first order in derivatives because the hopping Hamiltonian is quadratic.

It is convenient to write the hopping Hamiltonian in matrix form,
\[
\hat H_J = \sum_{j,k} h_{jk}\,\hat a_j^\dagger \hat a_k,
\]
where
\[
h_{jk}=
\begin{cases}
-J, & j,k \text{ nearest neighbors},\\
0, & \text{otherwise}.
\end{cases}
\]
Then the coherent drift generated by hopping is
\[
\frac{\partial P}{\partial t}\Big|_{J}
=
-\sum_j \frac{\partial}{\partial \alpha_j}
\left(
-i\sum_k h_{jk}\alpha_k\,P
\right)
-\sum_j \frac{\partial}{\partial \widetilde{\alpha}_j}
\left(
+i\sum_k \widetilde{\alpha}_k h_{kj}\,P
\right).
\]
For real symmetric nearest-neighbor hopping, $h_{jk}=h_{kj}$, this becomes
\[
\frac{\partial P}{\partial t}\Big|_{J}
=
-\sum_j \frac{\partial}{\partial \alpha_j}
\left(
+iJ\sum_{k\in \mathrm{nn}(j)} \alpha_k\,P
\right)
-\sum_j \frac{\partial}{\partial \widetilde{\alpha}_j}
\left(
-iJ\sum_{k\in \mathrm{nn}(j)} \widetilde{\alpha}_k\,P
\right).
\]
So hopping transports amplitudes between neighboring sites but does not introduce any new noise.

\section*{Fokker-Planck Equation}
Collecting the interaction, dissipation, and hopping pieces, the positive-$P$ distribution obeys a Fokker-Planck equation of the form
\[
\frac{\partial P}{\partial t}
=
-\sum_\mu \frac{\partial}{\partial x_\mu}\left(A_\mu P\right)
+ \frac12 \sum_{\mu,\nu}
\frac{\partial^2}{\partial x_\mu \partial x_\nu}
\left(D_{\mu\nu}P\right),
\]
with phase-space variables
\[
\bm{x}
=
(\alpha_1,\widetilde{\alpha}_1,\dots,\alpha_M,\widetilde{\alpha}_M)^T.
\]

For each site $j$, the drift components are
\begin{align}
A_{\alpha_j}
&=
-\frac{\gamma}{2}\alpha_j
-iU\,\widetilde{\alpha}_j\alpha_j^2
-i\sum_k h_{jk}\alpha_k, \\
A_{\widetilde{\alpha}_j}
&=
-\frac{\gamma}{2}\widetilde{\alpha}_j
+iU\,\widetilde{\alpha}_j^{\,2}\alpha_j
+i\sum_k \widetilde{\alpha}_k h_{kj}.
\end{align}
For real nearest-neighbor hopping this is
\begin{align}
A_{\alpha_j}
&=
-\frac{\gamma}{2}\alpha_j
-iU\,\widetilde{\alpha}_j\alpha_j^2
+iJ\sum_{k\in \mathrm{nn}(j)} \alpha_k, \\
A_{\widetilde{\alpha}_j}
&=
-\frac{\gamma}{2}\widetilde{\alpha}_j
+iU\,\widetilde{\alpha}_j^{\,2}\alpha_j
-iJ\sum_{k\in \mathrm{nn}(j)} \widetilde{\alpha}_k.
\end{align}

The diffusion is block diagonal in the site index:
\[
D^{(j)}
=
\begin{pmatrix}
-iU\,\alpha_j^2 & 0 \\
0 & +iU\,\widetilde{\alpha}_j^{\,2}
\end{pmatrix}.
\]
One convenient factorization is
\[
B^{(j)}
=
\begin{pmatrix}
\sqrt{-iU}\,\alpha_j & 0 \\
0 & \sqrt{+iU}\,\widetilde{\alpha}_j
\end{pmatrix},
\qquad
D^{(j)} = B^{(j)} \left(B^{(j)}\right)^T.
\]
This choice is especially useful because it leads directly to the same real-noise form used in the code.

\section*{Ito Stochastic Differential Equations}
The Fokker-Planck equation is equivalent to Ito SDEs,
\[
d\bm{x} = \bm{A}\,dt + B\,d\bm{W},
\]
where the Wiener increments satisfy
\[
\langle dW_\mu \rangle = 0,
\qquad
\langle dW_\mu\,dW_\nu \rangle = \delta_{\mu\nu}\,dt.
\]
Using the factorization above, the site-resolved positive-$P$ equations become
\begin{align}
d\alpha_j
&=
\left(
-\frac{\gamma}{2}\alpha_j
-iU\,\widetilde{\alpha}_j\alpha_j^2
-i\sum_k h_{jk}\alpha_k
\right)dt
+ \sqrt{-iU}\,\alpha_j\,dW_{j,1}, \\
d\widetilde{\alpha}_j
&=
\left(
-\frac{\gamma}{2}\widetilde{\alpha}_j
+iU\,\widetilde{\alpha}_j^{\,2}\alpha_j
+i\sum_k \widetilde{\alpha}_k h_{kj}
\right)dt
+ \sqrt{+iU}\,\widetilde{\alpha}_j\,dW_{j,2}.
\end{align}
For real nearest-neighbor hopping,
\begin{align}
d\alpha_j
&=
\left(
-\frac{\gamma}{2}\alpha_j
-iU\,\widetilde{\alpha}_j\alpha_j^2
+iJ\sum_{k\in \mathrm{nn}(j)} \alpha_k
\right)dt
+ \sqrt{-iU}\,\alpha_j\,dW_{j,1}, \\
d\widetilde{\alpha}_j
&=
\left(
-\frac{\gamma}{2}\widetilde{\alpha}_j
+iU\,\widetilde{\alpha}_j^{\,2}\alpha_j
-iJ\sum_{k\in \mathrm{nn}(j)} \widetilde{\alpha}_k
\right)dt
+ \sqrt{+iU}\,\widetilde{\alpha}_j\,dW_{j,2}.
\end{align}
The noises $dW_{j,1}$ and $dW_{j,2}$ are independent real Wiener increments for each site.

\section*{Connection to Observables}
Normal-ordered observables are obtained by stochastic averaging. In particular,
\begin{align}
\langle \hat a_j \rangle &= \langle \alpha_j \rangle_{\mathrm{stoch}}, \\
\langle \hat a_j^\dagger \rangle &= \langle \widetilde{\alpha}_j \rangle_{\mathrm{stoch}}, \\
\langle \hat a_j^\dagger \hat a_j \rangle
&=
\langle \widetilde{\alpha}_j \alpha_j \rangle_{\mathrm{stoch}}, \\
\langle \hat a_j^{\dagger 2}\hat a_j^2 \rangle
&=
\langle (\widetilde{\alpha}_j \alpha_j)^2 \rangle_{\mathrm{stoch}}.
\end{align}
This is the reason the code estimates the mean occupation from $\widetilde{\alpha}\alpha$ and the factorial second moment from $(\widetilde{\alpha}\alpha)^2$.

\section*{Specialization to the Current Code}
The file \texttt{src/bosonic\_dissipation/positive\_p\_method.py} currently implements the \emph{single-site} interacting case. Therefore $M=1$ and there are no neighbors, so the hopping sums vanish:
\[
\sum_{k\in \mathrm{nn}(1)} \alpha_k = 0,
\qquad
\sum_{k\in \mathrm{nn}(1)} \widetilde{\alpha}_k = 0.
\]
The general lattice SDEs reduce to
\begin{align}
d\alpha
&=
\left(
-\frac{\gamma}{2}\alpha
-iU\,\widetilde{\alpha}\alpha^2
\right)dt
+ \sqrt{-iU}\,\alpha\,dW_1, \\
d\widetilde{\alpha}
&=
\left(
-\frac{\gamma}{2}\widetilde{\alpha}
+iU\,\widetilde{\alpha}^{\,2}\alpha
\right)dt
+ \sqrt{+iU}\,\widetilde{\alpha}\,dW_2.
\end{align}
These are exactly the SDEs used in the solver, with the code identifications
\[
U \leftrightarrow \texttt{interaction\_strength},
\qquad
\widetilde{\alpha} \leftrightarrow \texttt{alpha\_plus}.
\]
In Euler-Maruyama form, the code updates are
\begin{align}
\alpha(t+\Delta t)
&=
\alpha(t)
+
\left(
-\frac{\gamma}{2}\alpha
-iU\,\widetilde{\alpha}\alpha^2
\right)\Delta t
+ \sqrt{-iU}\,\alpha\,\Delta W_1, \\
\widetilde{\alpha}(t+\Delta t)
&=
\widetilde{\alpha}(t)
+
\left(
-\frac{\gamma}{2}\widetilde{\alpha}
+iU\,\widetilde{\alpha}^{\,2}\alpha
\right)\Delta t
+ \sqrt{+iU}\,\widetilde{\alpha}\,\Delta W_2,
\end{align}
with independent Gaussian increments
\[
\Delta W_1,\Delta W_2 \sim \mathcal{N}(0,\Delta t).
\]
This matches the implementation
\begin{align*}
\texttt{drift\_alpha} &= -0.5 * \gamma * \alpha - 1j * U * \alpha * \widetilde{\alpha} * \alpha, \\
\texttt{drift\_alpha\_plus} &= -0.5 * \gamma * \widetilde{\alpha} + 1j * U * \widetilde{\alpha} * \alpha * \widetilde{\alpha}, \\
\alpha &\leftarrow \alpha + \texttt{drift\_alpha} * dt + \sqrt{-iU}\,\alpha\,\texttt{noise\_1}, \\
\widetilde{\alpha} &\leftarrow \widetilde{\alpha} + \texttt{drift\_alpha\_plus} * dt + \sqrt{+iU}\,\widetilde{\alpha}\,\texttt{noise\_2},
\end{align*}
where \texttt{noise\_1} and \texttt{noise\_2} are sampled as independent real Gaussians with variance $dt$.

\section*{Summary}
For the open Bose-Hubbard model, the positive-$P$ method turns the Lindblad master equation into stochastic equations for independent complex variables $\alpha_j$ and $\widetilde{\alpha}_j$. The interaction $U$ generates both nonlinear drift and multiplicative noise, the dissipation $\gamma$ contributes linear damping, and the hopping $J$ produces deterministic coupling between neighboring sites. The resulting Ito SDEs are
\begin{align*}
d\alpha_j
&=
\left(
-\frac{\gamma}{2}\alpha_j
-iU\,\widetilde{\alpha}_j\alpha_j^2
+iJ\sum_{k\in \mathrm{nn}(j)} \alpha_k
\right)dt
+ \sqrt{-iU}\,\alpha_j\,dW_{j,1}, \\
d\widetilde{\alpha}_j
&=
\left(
-\frac{\gamma}{2}\widetilde{\alpha}_j
+iU\,\widetilde{\alpha}_j^{\,2}\alpha_j
-iJ\sum_{k\in \mathrm{nn}(j)} \widetilde{\alpha}_k
\right)dt
+ \sqrt{+iU}\,\widetilde{\alpha}_j\,dW_{j,2}.
\end{align*}
For the current single-site code, the $J$ terms vanish and the remaining equations match the implemented solver exactly.

\end{document}
